\documentclass[spanish, a4paper,twoside]{xc_report}
\usepackage[spanish]{babel}
\usepackage[T1]{fontenc}
\usepackage[utf8]{inputenc}
\usepackage{pdfpages}
\usepackage[font=small]{caption}
\usepackage{multicol}
\usepackage{subcaption}

\author{XC ingeniería estructural}
\date{\today}
\renewcommand{\revision}{0.0}
\title{Cálculo sísmico}
\setcounter{tocdepth}{3}
\begin{document}

\newenvironment{Figure}
  {\par\medskip\noindent\minipage{\linewidth}}
  {\endminipage\par\medskip}
  
Se emplean, como en el proyecto base, acelerogramas generados artificialmente, que resultan compatibles con el espectro definido en UNE-EN 1998-1\S3.3.1.4.2 en el emplazamiento de la estructura, para un amortiguamiento del 5\%.

El procediemiento empleado es el siguiente:
\begin{enumerate}
\item Obtención espectro de respuesta elástico acorde al anejo nacional de UNE EN 1998-1
\item Conforme a UNE-EN 1998-2\S3.2.3(1)P se generan artificialmente tres pares de acelerogramas de cada tipo; en el apéndice \ref{scc:accel_spectr_horiz} se pueden consultar los correspondientes a la aceleración horizontal y en el \ref{scc:accel_spectr_vert} los de la vertical.
\item Para garantizar la consistencia de estos acelerogramas con el espectro de respuesta elásticos se escala la amplitud del movimiento empleando el siguiente procedimiento:
  \begin{enumerate}
    \item Para cada par de acelerogramas se determina el espectro que corresponde a cada uno de ellos y el espectro SRSS del par tomando, como aceleración de éste, la raíz cuadrada de la suma de los cuadrados (SRSS) de las aceleraciones del espectro de cada componente (ver apéndice \ref{scc:spect_SRSS})
    
  \end{enumerate}
  

\end{enumerate}

%\portadilla{Apéndices.}
\clearpage
\appendix

\section{Acelerogramas sintéticos y espectros de respuesta horizontal}\label{scc:accel_spectr_horiz}
\begin{figure}[ht]
  \centering
  \begin{subfigure}[t]{0.42\textwidth}
    \centering
    \includegraphics[width=\textwidth]{output/horiz_accel/motion1/motion1_timehist.png}
    \caption{Aceleraciones, velocidades y desplazamientos horizontales.}
  \end{subfigure}
  \hfill
  \begin{subfigure}[t]{0.42\textwidth}
      \centering
      \includegraphics[width=\textwidth]{output/horiz_accel/motion1/motion1_spectra.png}
      \caption{Espectro de respuesta}
      \end{subfigure}

      \caption{Serie 1 de aceleración horizontal sintética}\label{fg:horiz_motion1}
\end{figure}

\begin{figure}[ht]
  \centering
  \begin{subfigure}[t]{0.42\textwidth}
    \centering
    \includegraphics[width=\textwidth]{output/horiz_accel/motion2/motion2_timehist.png}
    \caption{Aceleraciones, velocidades y desplazamientos horizontales.}
  \end{subfigure}
  \hfill
  \begin{subfigure}[t]{0.42\textwidth}
      \centering
      \includegraphics[width=\textwidth]{output/horiz_accel/motion2/motion2_spectra.png}
      \caption{Espectro de respuesta}
      \end{subfigure}

      \caption{Serie 2 de aceleración horizontal sintética}\label{fg:horiz_motion2}
\end{figure}

 
\begin{figure}[ht]
  \centering
  \begin{subfigure}[t]{0.42\textwidth}
    \centering
    \includegraphics[width=\textwidth]{output/horiz_accel/motion3/motion3_timehist.png}
    \caption{Aceleraciones, velocidades y desplazamientos horizontales.}
  \end{subfigure}
  \hfill
  \begin{subfigure}[t]{0.42\textwidth}
      \centering
      \includegraphics[width=\textwidth]{output/horiz_accel/motion3/motion3_spectra.png}
      \caption{Espectro de respuesta}
      \end{subfigure}

      \caption{Serie 3 de aceleración horizontal sintética}\label{fg:horiz_motion3}
\end{figure}

\begin{figure}[ht]
  \centering
  \begin{subfigure}[t]{0.42\textwidth}
    \centering
    \includegraphics[width=\textwidth]{output/horiz_accel/motion4/motion4_timehist.png}
    \caption{Aceleraciones, velocidades y desplazamientos horizontales.}
  \end{subfigure}
  \hfill
  \begin{subfigure}[t]{0.42\textwidth}
      \centering
      \includegraphics[width=\textwidth]{output/horiz_accel/motion4/motion4_spectra.png}
      \caption{Espectro de respuesta}
      \end{subfigure}

      \caption{Serie 4 de aceleración horizontal sintética}\label{fg:horiz_motion4}
\end{figure}

\begin{figure}[ht]
  \centering
  \begin{subfigure}[t]{0.42\textwidth}
    \centering
    \includegraphics[width=\textwidth]{output/horiz_accel/motion5/motion5_timehist.png}
    \caption{Aceleraciones, velocidades y desplazamientos horizontales.}
  \end{subfigure}
  \hfill
  \begin{subfigure}[t]{0.42\textwidth}
      \centering
      \includegraphics[width=\textwidth]{output/horiz_accel/motion5/motion5_spectra.png}
      \caption{Espectro de respuesta}
      \end{subfigure}

      \caption{Serie 5 de aceleración horizontal sintética}\label{fg:horiz_motion5}
\end{figure}

\begin{figure}[ht]
  \centering
  \begin{subfigure}[t]{0.42\textwidth}
    \centering
    \includegraphics[width=\textwidth]{output/horiz_accel/motion6/motion6_timehist.png}
    \caption{Aceleraciones, velocidades y desplazamientos horizontales.}
  \end{subfigure}
  \hfill
  \begin{subfigure}[t]{0.42\textwidth}
      \centering
      \includegraphics[width=\textwidth]{output/horiz_accel/motion6/motion6_spectra.png}
      \caption{Espectro de respuesta}
      \end{subfigure}

      \caption{Serie 6 de aceleración horizontal sintética}\label{fg:horiz_motion6}
\end{figure}

\clearpage
\section{Acelerogramas sintéticos y espectros de respuesta vertical}\label{scc:accel_spectr_vert}
\begin{figure}[ht]
  \centering
  \begin{subfigure}[t]{0.42\textwidth}
    \centering
    \includegraphics[width=\textwidth]{output/vertical_accel/motion1/motion1_timehist.png}
    \caption{Aceleraciones, velocidades y desplazamientos verticales.}
  \end{subfigure}
  \hfill
  \begin{subfigure}[t]{0.42\textwidth}
      \centering
      \includegraphics[width=\textwidth]{output/vertical_accel/motion1/motion1_spectra.png}
      \caption{Espectro de respuesta}
      \end{subfigure}

      \caption{Serie 1 de aceleración vertical sintética}\label{fg:vert_motion1}
\end{figure}

\begin{figure}[ht]
  \centering
  \begin{subfigure}[t]{0.42\textwidth}
    \centering
    \includegraphics[width=\textwidth]{output/vertical_accel/motion2/motion2_timehist.png}
    \caption{Aceleraciones, velocidades y desplazamientos verticales.}
  \end{subfigure}
  \hfill
  \begin{subfigure}[t]{0.42\textwidth}
      \centering
      \includegraphics[width=\textwidth]{output/vertical_accel/motion2/motion2_spectra.png}
      \caption{Espectro de respuesta}
      \end{subfigure}

      \caption{Serie 2 de aceleración vertical sintética}\label{fg:vert_motion2}
\end{figure}

 
\begin{figure}[ht]
  \centering
  \begin{subfigure}[t]{0.42\textwidth}
    \centering
    \includegraphics[width=\textwidth]{output/vertical_accel/motion3/motion3_timehist.png}
    \caption{Aceleraciones, velocidades y desplazamientos verticales.}
  \end{subfigure}
  \hfill
  \begin{subfigure}[t]{0.42\textwidth}
      \centering
      \includegraphics[width=\textwidth]{output/vertical_accel/motion3/motion3_spectra.png}
      \caption{Espectro de respuesta}
      \end{subfigure}

      \caption{Serie 3 de aceleración vertical sintética}\label{fg:vert_motion3}
\end{figure}

\begin{figure}[ht]
  \centering
  \begin{subfigure}[t]{0.42\textwidth}
    \centering
    \includegraphics[width=\textwidth]{output/vertical_accel/motion4/motion4_timehist.png}
    \caption{Aceleraciones, velocidades y desplazamientos verticales.}
  \end{subfigure}
  \hfill
  \begin{subfigure}[t]{0.42\textwidth}
      \centering
      \includegraphics[width=\textwidth]{output/vertical_accel/motion4/motion4_spectra.png}
      \caption{Espectro de respuesta}
      \end{subfigure}

      \caption{Serie 4 de aceleración vertical sintética}\label{fg:vert_motion4}
\end{figure}

\begin{figure}[ht]
  \centering
  \begin{subfigure}[t]{0.42\textwidth}
    \centering
    \includegraphics[width=\textwidth]{output/vertical_accel/motion5/motion5_timehist.png}
    \caption{Aceleraciones, velocidades y desplazamientos verticales.}
  \end{subfigure}
  \hfill
  \begin{subfigure}[t]{0.42\textwidth}
      \centering
      \includegraphics[width=\textwidth]{output/vertical_accel/motion5/motion5_spectra.png}
      \caption{Espectro de respuesta}
      \end{subfigure}

      \caption{Serie 5 de aceleración vertical sintética}\label{fg:vert_motion5}
\end{figure}

\begin{figure}[ht]
  \centering
  \begin{subfigure}[t]{0.42\textwidth}
    \centering
    \includegraphics[width=\textwidth]{output/vertical_accel/motion6/motion6_timehist.png}
    \caption{Aceleraciones, velocidades y desplazamientos verticales.}
  \end{subfigure}
  \hfill
  \begin{subfigure}[t]{0.42\textwidth}
      \centering
      \includegraphics[width=\textwidth]{output/vertical_accel/motion6/motion6_spectra.png}
      \caption{Espectro de respuesta}
      \end{subfigure}

      \caption{Serie 6 de aceleración vertical sintética}\label{fg:vert_motion6}
\end{figure}

\clearpage
\begin{multicols}{2}
  \section{Espectros SRSS}\label{scc:spect_SRSS}
  % Horizontales
\begin{Figure}
  \centering
  \includegraphics[width=\linewidth]{output/horiz_accel/figures/spectr_pair1.png}
  \captionof{figure}{Espectros de respuesta del par de acelerogramas horizontales 1 y 2 y espectro SRSS} \label{fg:horiz_spect_SRSS_pair1}
\end{Figure}
\begin{Figure}
  \centering
  \includegraphics[width=\linewidth]{output/horiz_accel/figures/spectr_pair2.png}
  \captionof{figure}{Espectros de respuesta del par de acelerogramas horizontales 3 y 4 y espectro SRSS} \label{fg:horiz_spect_SRSS_pair2}
\end{Figure}
\begin{Figure}
  \centering
  \includegraphics[width=\linewidth]{output/horiz_accel/figures/spectr_pair3.png}
  \captionof{figure}{Espectros de respuesta del par de acelerogramas horizontales 5 y 6 y espectro SRSS} \label{fg:horiz_spect_SRSS_pair3}
\end{Figure}

% Verticales
\begin{Figure}
  \centering
  \includegraphics[width=\linewidth]{output/vertical_accel/figures/spectr_pair1.png}
  \captionof{figure}{Espectros de respuesta del par de acelerogramas verticales 1 y 2 y espectro SRSS} \label{fg:vert_spect_SRSS_pair1}
\end{Figure}
\begin{Figure}
  \centering
  \includegraphics[width=\linewidth]{output/vertical_accel/figures/spectr_pair2.png}
  \captionof{figure}{Espectros de respuesta del par de acelerogramas verticales 3 y 4 y espectro SRSS} \label{fg:vert_spect_SRSS_pair2}
\end{Figure}
\begin{Figure}
  \centering
  \includegraphics[width=\linewidth]{output/vertical_accel/figures/spectr_pair3.png}
  \captionof{figure}{Espectros de respuesta del par de acelerogramas verticales 5 y 6 y espectro SRSS} \label{fg:vert_spect_SRSS_pair3}
\end{Figure}

\end{multicols}
\end{document}

